\section{The integral lattices of the theta lifts}\label{sec:integral-lifts}

Keep $T,u_0,u,p$ as in \S\ref{sec:chain-divisibility}, with $p$ unramified
in $E$. We check the two local assertions used in
Theorem~\ref{thm:chain-divisibility}. Fix a lift $j$ and an orbit $v$.
Write $n=N_{j-1}$, $d=N_j$ and $w=w_j(\tau_s)$ at position $s$.

\subsection{The edge and the two polynomial transfers}

For a dominant weight $\beta$ of length $d$, put
$\bar\beta=\beta-\beta_d(1,\ldots,1)$. In the standard frame take the
$\mathcal O$-span of the products $\prod_a\Delta_{S_a}(Y)$, where
$\Delta_S(Y)=\det Y_{S,[|S|]}$ and $|S_a|$ are the column heights
of $\bar\beta$. Include the determinant twist of
\citep[Definition~11.3]{adriazola-grh}. Changes of integral basis preserve
the span by Cauchy--Binet in both directions. In an adapted basis an
edge of type $i$ has matrix $\operatorname{diag}(pI_i,I_{d-i})$.
Expressing such a product from the smaller lattice in the larger frame
gives the factor
\begin{equation}\label{eq:true-edge-factor}
                 p^{-i\beta_d-\sum_a|S_a\cap[i]|}.
\end{equation}
The twist contributes $p^{-i\beta_d}$ and a row in $[i]$
contributes $p^{-1}$ at each occurrence. Since
$\sum_a\min(|S_a|,i)=\sum_{\ell\le i}(\beta_\ell-\beta_d)$,
the sum of the first $i$ entries of $\beta$ clears the denominator.
At level zero this gives $y_v I_0\subseteq I_0$ with the exponent $A_v$.

Write $\lambda=\lambda^{(j-1)}_{\tau_s}$ and let
$h_1,\ldots,h_c$ be the column heights of the partition
$\mu=(\lambda_1+w,\ldots,\lambda_n+w,0,\ldots,0)$ of length $d$.
Thus $c=\lambda_1+w$ and $|\mu|=\sum_a h_a$. Let $D$ be an
$n\times d$ matrix of differential symbols and $Y$ a $d\times d$
target frame. Complements of index sets use the dimension of the
corresponding row or column space.

\begin{proposition}\label{prop:true-polynomials}
On a generator of the source lattice, each term of the even transfer
has the form
\[
          \prod_a\det D_{R_a,I_a}\ \prod_a\Delta_{I_a}(Y),
               \qquad |R_a|=|I_a|=h_a .
\]
Each term of the odd adjugate transfer has the form
\[
       \pm\prod_a\det D_{R_a^c,I_a}\ \prod_a\Delta_{I_a^c}(Y),
               \qquad |R_a|=n-h_a,\quad |I_a|=h_a .
\]
With the normalization $p^{n^{(j)}_{v,i}}$, the edge exponents
are those of \S\ref{sec:chain-divisibility} at both parities.
\end{proposition}

\begin{proof}
At an even position the transfer substitutes the first $n$ columns of
$DY$ into a source generator $G$, and multiplies by the determinant
of that block to the power $\lambda_n+w\ge0$. These extra columns
have height $n$. Apply Cauchy--Binet to each column of $\mu$.
The target bottom weight is $-w$, and the even normalization is
$\sum_{\ell\le i}\lambda_\ell=-iw+\sum_a\min(h_a,i)$.
Formula~\eqref{eq:true-edge-factor} leaves
$\sum_a(\min(h_a,i)-|I_a\cap[i]|)$.

At an odd position the source and target use the dual weights.
Their column heights are respectively $n-h_a$ and $d-h_a$ in reverse
order, allowing empty columns. Write $J_h$ for the last $h$ indices
in the ambient row set. The adjugate substitution is \citep{alpoge-local-symbol}
\[
 U=(\operatorname{adj}Y)_{J_n,:}D^{\mathsf T},\qquad
 (\det Y)^{c-|\mu|}(\det U)^{c-nc+|\mu|}G(\operatorname{adj}U).
\]
For $|R|=n-h$, Jacobi's identity at the two matrices gives
\[
 \Delta_R(\operatorname{adj}U)
 =\pm(\det U)^{n-h-1}\det U_{J_h,R^c},\qquad
 \det(\operatorname{adj}Y)_{J_h,I}
 =\pm(\det Y)^{h-1}\Delta_{I^c}(Y).
\]
Between these identities, Cauchy--Binet expands $\det U_{J_h,R^c}$
in the minors $\det D_{R^c,I}$.
The exponents sum to $nc-|\mu|-c$ for $\det U$ and $|\mu|-c$
for $\det Y$. Both cancel. The resulting polynomial identity over
$\mathbb Z$ includes the empty minors.
Now $\beta_d=-\lambda_1$ and the odd normalization is $iw$. By
\eqref{eq:true-edge-factor} the remaining exponent is
\[
 iw+i\lambda_1-\sum_a|I_a^c\cap[i]|=\sum_a|I_a\cap[i]| .\qedhere
\]
\end{proof}

\subsection{What the shared frame detects}

For a normalized source partition $\nu$ of either parity put
$H_\nu(X)=\prod_{h}\Delta_{[h]}(X)$, taking the heights of its
columns with multiplicity. Let $C_\nu$ be the span of the
coefficients in $Z$ of $H_\nu(ZX)$ for one matrix $Z$ of indeterminates.

\begin{lemma}\label{lem:true-common-frame}
At either parity the corrected shared-frame bounds give the transfer
bound on $C_\nu$.
\end{lemma}

\begin{proof}
At an even position apply the transfer to $H_\nu(ZX)$ and put $k_0=\lambda_n+w$.
After multiplying by $(\det Z)^{k_0}$, the coefficient attached in the
Cauchy--Binet expansion to a column choice $(I_a)_a$ is the pairing with
\[
             \prod_a\det (ZD)_{[h_a],I_a}.
\]
The $k_0$ additional full-height minors each supply one $\det Z$.

At an odd position $Z\operatorname{adj}U=(\det Z)\operatorname{adj}(UZ^{-1})$.
For a column of height $n-h$, Jacobi's identity therefore gives
\[
 \Delta_{[n-h]}(Z\operatorname{adj}U)
    =(\det Z)(\det U)^{n-h-1}\det (UZ^{-1})_{J_h,J_h}.
\]
Let $J$ reverse the $n$ rows. Expand as in Proposition~\ref{prop:true-polynomials},
now with $Z^{-\mathsf T}D$ in place of $D$. Multiplying the transfer
by $(\det Z)^{|\mu|-c}$ gives, up to signs, the shared-frame pairings
with $Z$ replaced by $J(\operatorname{adj}Z)^{\mathsf T}$.
This substitution is integral. At either parity the edge correction depends
only on $(I_a)_a$. A power of $\det Z$ preserves the minimum coefficient
valuation by Gauss's lemma. Taking coefficients in $Z$ proves the claim.
\end{proof}

The full lattice can be larger. For a single row of length $p$ and
two source indices, take $f(D)=D_{11}^{p-1}D_{21}$. With every column
index equal to $1$ the shared-frame jet is
\[
 (z_1\partial_{11}+z_2\partial_{21})^p f(0)
     =p!\,z_1^{p-1}z_2,\qquad
 \partial_{11}^{p-1}\partial_{21} f(0)=(p-1)! .
\]
Every choice involving another column vanishes. Thus divisibility of
all shared-frame coefficients by $p$ does not imply divisibility of
each individual row pairing in Proposition~\ref{prop:true-polynomials}.

The determinant factors can also be computed on the symbols. Put
$Y'=(g')^{\mathsf T}Y$ and $(gG)(X)=G(g^{\mathsf T}X)$. At an even
position take $D'=g^{\mathsf T}D(g')^{-\mathsf T}$. Then
$D'Y'=g^{\mathsf T}DY$ and the extra determinant gives
$(\det g)^{\lambda_n+w}$ times the transfer of $gG$.
At an odd position take $D'=g^{-1}Dg'$. The matrix $U$ in the proof
of Proposition~\ref{prop:true-polynomials} changes to
$(\det g')U g^{-\mathsf T}$. Since $\deg G=nc-|\mu|$, substitution
gives $(\det g'/\det g)^c$ times the transfer of $gG$.
Including the source and target twists gives, respectively,
\[
 (\det g)^{-\lambda_n}(\det g)^{\lambda_n+w}(\det g')^{-w}
      =(\det g/\det g')^w,\qquad
 (\det g/\det g')^{\lambda_1-c}=(\det g/\det g')^{-w}.
\]
The series uses the inverse transpose of the substitution on symbols,
by the chain rule for $P(\partial)f(0)$.

\begin{hypothesis}[\openinput]\label{hyp:true-jet-model}
On one local coefficient model for the prescribed lifts at $T,u_0$,
the polynomial transfers above give the jets in $A_j$ with these
changes of frame, the length factor of \citep[Definition~11.9]{adriazola-grh},
and integral factors at the remaining positions.
For each prescribed kernel, the corrected bounds on the shared-frame
coefficients in Hypothesis~\ref{hyp:chain-lattice-transfer} extend to
the transfer on the full source lattice in the adapted frames.
\end{hypothesis}

One sufficient check for the last assertion is
\[
 v_p\!\left(\left\langle\prod_{v,s,a}\det D_{v,s;R_a,I_a},f\right\rangle\right)
                     +\sum_v E_v\ge b
\]
for the row choices in Proposition~\ref{prop:true-polynomials}, using
the complementary source rows at odd positions. Here $f$ is the
prescribed kernel and $b$ is the shared-frame bound.

\subsection{The split Hecke relation}

Write $\ell(L)=-v_p(\det g)$ at $w\mid v$ and $q=q_v$.
The length exponent is the same at every lift,
\[
 \gamma_v=\sum_{s\ {\rm even}}w_j(\tau_s)-\sum_{s\ {\rm odd}}w_j(\tau_s)
          =\sum_{s\ {\rm even}}b_1(\tau_s)-\sum_{s\ {\rm odd}}b_1(\tau_s).
\]
Indeed \eqref{eq:chain-nodes} adds the same constant at the
$\nu_v$ positions of each parity.
Remove $p^{\gamma_v(\ell(L)-\ell(L'))}$ from the local step map as in
\citep[Definition~11.9]{adriazola-grh}, and write $s_j(L,L')$ for
the remaining frame scalar multiplying the character-normalized theta series.

\begin{hypothesis}[\openinput]\label{hyp:true-frame-scalar}
For the word $u$, the level data away from $v$ stay fixed under the
edges, with the unramified lattice indicator at $v$. In the coefficient
model of the lift, the adjusted frame scalar satisfies, for all moves at $v$,
\[
 \frac{s_j(L,L')}{s_j(L_0,L'_0)}
       =a^{\ell(L)-\ell(L_0)}(a')^{\ell(L')-\ell(L'_0)}
       \qquad(a,a'\in\mathcal O^\times).
\]
The units depend on $v,j,u$ and are common to the class sum.
\end{hypothesis}

The determinant-idele model of \citep[Definition~11.23 and
Proposition~11.27]{adriazola-grh} gives this ratio explicitly.
Let $\varpi_v$ have components $(p,p^{-1})$ at $w,\bar w$ and one
elsewhere. Its step character $\psi_j$, unramified at $v$ and of
type $w_j(\tau)+c_j$ on norm-one elements, satisfies
$v_p(\psi_j(\varpi_v))=\gamma_v$ in the fixed coefficient embedding.
Indeed if the ideal $(\varpi_v)^h=(\beta)$, then
$(\varpi_v)^{2h}=(\beta/\bar\beta)$. The quotient idele is a unit
everywhere, and its character value has valuation zero. Evaluating the
type on $\beta/\bar\beta$ gives $2h\gamma_v$.
With $\kappa=\chi((\varpi_v))$ the unit from the determinant character
and $t_\ell=\ell(t_r+1)/2$, the model gives
\[
 a=p^{-\gamma_v}\psi_j(\varpi_v)\kappa^{t_{j-1}},\qquad
 a'=p^{\gamma_v}\psi_j(\varpi_v)^{-1}\kappa^{-t_j}.
\]
Changing the two determinant ideles by $\varpi_v^{-\Delta\ell}$
and $\varpi_v^{-\Delta\ell'}$ up to local units gives the stated ratio.
Hypothesis~\ref{hyp:true-frame-scalar} requires it for the frames
used in \eqref{eq:chain-class-sum}.

\begin{proposition}\label{prop:true-hecke}
Under Hypothesis~\ref{hyp:true-frame-scalar}, the lift satisfies
\eqref{eq:chain-one-hecke} with
$\epsilon_{vj}=(aa')^{-1}$ and $\eta_{vj}=p^{\gamma_v}a$.
\end{proposition}

\begin{proof}
Put $g=d-n=t_r+1$ and $Q=q^2$. For the split theta kernel the normalized
series $\operatorname{vol}(\operatorname{Hom}(L,L'))^{1/2}
\chi_F^{-1}\theta_{L,L'}$ depends linearly on the indicator
$\mathbf1_{\{xL\subseteq L'\}}$ when the other data are fixed
\citep[Proposition~11.20]{adriazola-grh}. The volume is
$q^{2(n\ell'-d\ell)}$ by change of basis. Thus the local step map is
a fixed linear function of this indicator multiplied by
\[
      p^{\gamma_v(\ell-\ell')}a^\ell(a')^{\ell'}q^{d\ell-n\ell'},
      \qquad \ell=\ell(L),\quad \ell'=\ell(L').
\]
Shrinking the target by $i$ multiplies this scalar by
$p^{\gamma_vi}(a')^{-i}q^{ni}$. Enlarging the source by $k$ multiplies
it by $p^{\gamma_vk}a^kq^{dk}$.

If $xL\subseteq L'$, write $s$ for the rank of its reduction modulo
$p$. The target sublattices of type $i$ containing $xL$ are counted by
$\genfrac{[}{]}{0pt}{}{d-s}{i}_Q$, where the brackets count subspaces
over $\mathbb F_Q$. The source overlattices of type $k$ mapped into $L'$
are counted by $\genfrac{[}{]}{0pt}{}{n-s}{k}_Q$. Both counts vanish
if $xL\not\subseteq L'$. Splitting the subspace count by intersection
dimension gives
\[
 \genfrac{[}{]}{0pt}{}{d-s}{i}_Q=\sum_{l=0}^{\min(g,i)}
 Q^{(i-l)(g-l)}\genfrac{[}{]}{0pt}{}{g}{l}_Q
 \genfrac{[}{]}{0pt}{}{n-s}{i-l}_Q .
\]
In the class sum one may sum the source edges from either end. Each
edge orbit has weight the inverse order of its stabilizer.
Insert the two scalar ratios with $k=i-l$ and multiply by
$q^{-i(d-i)}$. Relative to the normalized source operator of degree
$i-l$, the coefficient is
\[
 (aa')^{-i}(p^{\gamma_v}a)^l
 q^{-l(g-l)}\genfrac{[}{]}{0pt}{}{g}{l}_{q^2}.
\]
By the Gaussian binomial theorem, the product following
$(p^{\gamma_v}a)^l$ is the elementary symmetric polynomial of degree
$l$ in $q^{g-1},q^{g-3},\ldots,q^{1-g}$. Summing in $i$ gives
\eqref{eq:chain-one-hecke} modulo $X^{r+1}$.
\end{proof}

\subsection{The prescribed lattices}

\begin{corollary}\label{cor:true-integral-chain}
Assume Hypotheses~\ref{hyp:inert-kernel-integrality},
\ref{hyp:true-jet-model} and~\ref{hyp:true-frame-scalar}. For the same
$T,u_0$ and each word $u$, $\mathcal K_m D_B\subseteq p^{G_B}I_m$.
\end{corollary}

\begin{proof}
In the frames of Hypothesis~\ref{hyp:true-jet-model} the determinant
ratios have total valuation $\sum_v\gamma_v(\ell(L''_v)-\ell(L_v))$.
This cancels the length factor of the step map. Proposition~\ref{prop:true-polynomials}
and the full-lattice assertion of Hypothesis~\ref{hyp:true-jet-model}
give \eqref{eq:chain-entry-inclusion}. Hypothesis~\ref{hyp:inert-kernel-integrality}
supplies its jet bounds by Corollary~\ref{cor:inert-chain-jets}.
Proposition~\ref{prop:true-hecke} supplies Hypothesis~\ref{hyp:chain-hecke}.
Apply Theorem~\ref{thm:chain-divisibility}.
\end{proof}

For the source comparison in Hypothesis~\ref{hyp:chain-endpoints},
consider a contributing $p\ne T_u$ and $M\subseteq I_0$. If
$S'_v\ge S_v$ for $v\in B$, the corresponding source module satisfies
\[
 P(M)=\sum_{C\subseteq B}y_C^m
        \bigl(p^{\sum_{v\in B\setminus C}S'_v}M\bigr)\subseteq D_B .
\]
The construction of \citep[Theorem~12.18 and Definition~12.20]{adriazola-grh}
takes an intersection $\Lambda\subseteq P(M)$ and then sets
$N_{u,\mathfrak p}=\pi^{-\lfloor e s_u(\mathfrak p)\rfloor}\Lambda$.
For this realization the source inclusion holds with
$a=\max(0,\lfloor e s_u(\mathfrak p)\rfloor/e)$.
The target inclusion is equivalent to
$v_p(P_u(F)(r_0))\ge-b$ on generators $F$ of $I_m$ and $r_0$ of
$R_{\mathfrak p}$. Identifying this source and controlling the two
losses for the prescribed period map remain the endpoint input
of Hypothesis~\ref{hyp:chain-endpoints}.
